\documentclass[11pt,a4paper]{article}

\usepackage[T1]{fontenc}
\usepackage[utf8]{inputenc}
\usepackage[spanish]{babel}
\usepackage{geometry}
\usepackage{microtype}
\usepackage{setspace}
\usepackage{booktabs}
\usepackage{tabularx}
\usepackage{enumitem}
\usepackage{hyperref}
\usepackage{xcolor}
\usepackage{titlesec}
\usepackage{parskip}

\geometry{margin=2.2cm}
\setstretch{1.12}
\setlist[itemize]{leftmargin=1.4em,itemsep=0.25em,topsep=0.35em}
\setlist[enumerate]{leftmargin=1.4em,itemsep=0.25em,topsep=0.35em}

\definecolor{upnavy}{HTML}{1F3A5F}
\hypersetup{colorlinks=true,urlcolor=upnavy,linkcolor=upnavy}
\titleformat{\section}{\large\bfseries\color{upnavy}}{\thesection.}{0.6em}{}
\titleformat{\subsection}{\normalsize\bfseries\color{upnavy}}{\thesubsection.}{0.5em}{}

\title{Cómo quedó ordenada la carpeta Minería}
\author{Johao Rafael Mendoza Fabián\\Universidad del Pacífico}
\date{Agosto 2026}

\begin{document}
\maketitle
\vspace{-1.2em}

\begin{center}
\textit{Nota para la reunión: qué había, cómo lo organicé y cómo se usa ahora.}
\end{center}

\section{En una frase}

La carpeta ya no es un depósito de archivos mezclados. Quedó armada como un archivo de investigación: primero se entiende el proyecto, después están los datos, el código, los resultados y, al final, lo que no hay que tocar para el trabajo del día a día.

Disco: \texttt{Johao\_asistente} $\rightarrow$ carpeta \textbf{Minería}.

\section{Qué problema había}

Todo estaba en la raíz o en carpetas con nombres poco claros. En el mismo sitio convivían la bitácora, JSON, Excel, Word, notebooks, papers, una segunda entrega y copias. Había duplicados (los PDF de la Defensoría, series NDVI/NDWI, el notebook de indicadores). Un compañero nuevo no sabía por dónde empezar ni qué archivo era el vigente.

\section{La lógica del orden}

Numeré las carpetas para que Drive las muestre de arriba hacia abajo, en el orden en que uno debería leerlas:

\begin{center}
\begin{tabular}{@{}clp{9.2cm}@{}}
\toprule
\textbf{Carpeta} & \textbf{Nombre} & \textbf{Para qué} \\
\midrule
00 & Inicio & De qué va el proyecto, bitácora, pendientes y reglas simples. \\
01 & Datos & Fuentes crudas por un lado; tablas y bases ya armadas por el otro. \\
02 & Código & Notebooks agrupados por tema (conflictos, satélite, mapas). \\
03 & Análisis & Mapas en HTML y fichas de minas. \\
04 & Literatura & Papers y material del seminario. \\
05 & Entregas & Lo que ya se mandó, sin mezclarlo con lo que se sigue usando. \\
99 & Archivo & Copias, código viejo y duplicados. \\
\bottomrule
\end{tabular}
\end{center}

\vspace{0.6em}
En la raíz de Minería solo quedan esas siete carpetas y un \texttt{LEEME.md}. Ese archivo es la puerta de entrada.

\section{Cómo se separan los datos}

En \texttt{01\_Datos} hay dos cajones a propósito:

\begin{itemize}
  \item \textbf{Crudos.} Lo que trajimos de afuera y no conviene editar: reportes de la Defensoría del Pueblo, JSON de OCMAL, shapefiles del IGN, PBI y demografía.
  \item \textbf{Procesados.} Lo que ya está listo para analizar: \texttt{conflictos\_data.json}, la BD de Defensoría, listados de minas e indicadores satelitales (NDVI, NDWI, NDMI).
\end{itemize}

Si alguien saca una tabla nueva, va a procesados. Los crudos se miran, no se pisan.

\section{Código y resultados}

El código quedó en \texttt{02\_Codigo}, por tema: extracción de conflictos, scraping de OCMAL, índices satelitales, mapas y segmentación de imágenes. Cada carpeta tiene el notebook que se usa ahora.

Los mapas para abrir en el navegador y las fichas (Toromocho, Quellaveco, Antamina) están en \texttt{03\_Analisis}, no mezclados con los PDF de papers. Las lecturas van a \texttt{04\_Literatura}. La segunda entrega (listado Huancayo e imágenes de esa fecha) está guardada en \texttt{05\_Entregas}, como snapshot.

\section{Qué hice con los duplicados}

No borré nada. Lo que estaba repetido ---la segunda pila de PDF de Defensoría, copias de Excel, el notebook viejo de indicadores, las series NDVI/NDWI duplicadas--- pasó a \texttt{99\_Archivo}, con nombres que dicen \texttt{DUPLICADO} cuando hace falta. El archivo vigente se busca en \texttt{01\_Datos} o \texttt{02\_Codigo}.

\section{Documentación para quien llegue después}

En cada carpeta grande hay un README corto, escrito como notas del equipo. En \texttt{00\_Inicio} está además:

\begin{itemize}
  \item cómo empezar, según si uno quiere datos, código o el estado del avance;
  \item convenciones de nombres y de dónde guardar cosas nuevas;
  \item la bitácora y los puntos pendientes.
\end{itemize}

En \texttt{01\_Datos} hay un diccionario: qué es cada JSON o Excel y para qué se usa. No hace falta conocer el historial del Drive para orientarse.

\section{Qué pedirle al equipo}

\begin{enumerate}
  \item Entrar por \texttt{LEEME.md} y \texttt{00\_Inicio}.
  \item No dejar archivos sueltos en la raíz.
  \item No editar \texttt{crudos}; las salidas nuevas van a \texttt{procesados} o a \texttt{03\_Analisis}.
  \item Las copias, a \texttt{99\_Archivo}, no al lado del archivo que sí se usa.
\end{enumerate}

Con eso la carpeta se mantiene sola. El contenido del proyecto no cambió: cambió el mapa para encontrarlo.

\end{document}
